\scanpage{178}
\mspage{94}
% archive sheet number: 177
% lines 1-3, continuation
\unsure{grde|grande} devant $\pi$ donc filtrante, cela implique qu'elle est $0$-connexe, et on peut appliquer le $=$ résultat à $\inserted{C'=}\prod_{j\in J'}C_j$, et on gagne.

% lines 4-12
Donc on est ramené à prouver que \deleted{\unsure{le}} $\prod C_j$ est grand devant $\pi$. Ceci \deleted{peut} se voit soit sur le critère $\mathrm{PF}_{\pi}1,\mathrm{PF}_{\pi}2$ de 4.4.4 b, soit \deleted{sur la définition} en se ramenant au cas \deleted{des} où les $C_j$ sont \unsure{ordonnées|ordonnés}, mais cela a l'air plus long. Notons que l'on prouve de $=$ qu'un \inserted{(petit)} produit de catégories filtrantes est encore \inserted{une cat.} filtrante.

% lines 13-20
\uline{Remarques} 1. Cet argument ne montre nullement ($=$ si $C$ est stable par $\unsure{\varinjlim|\underline{\lim}}$ de cardinal $\leq\pi$) que pour tout $\varphi\in\underline{\operatorname{Hom}}(J,\operatorname{Ind}(\deleted{\unsure{J|j}}C))$, le foncteur induit:
\[
\underline{\operatorname{Hom}}(\unsure{I|J},C)_{/\varphi}\longrightarrow C_{/\varphi(j)}
\tag*{(4.4.6.3)}
\]
est cofinal, et je doute qu'il le soit toujours. La difficulté provient du fait que $C\longrightarrow\operatorname{Ind}C$ ne commute pas aux sommes de cardinal $\leq\pi$, seulement aux sommes finies.
